\capitulo{3}{Conceptos teóricos}

Este capítulo presenta, de forma sencilla, los conceptos que sirven de base al sistema desarrollado. La idea principal del proyecto es que un paciente pueda escribir o dictar sus síntomas y que la aplicación proponga una especialidad médica adecuada. Para conseguirlo, el sistema transforma el texto del paciente en datos que un modelo de aprendizaje automático pueda interpretar y, a partir de ellos, realiza una clasificación.

No se pretende hacer una explicación matemática de cada algoritmo, sino justificar por qué se han utilizado estas técnicas y qué papel cumple cada una dentro del sistema. Todas las decisiones se han tomado buscando un equilibrio entre precisión, sencillez, rapidez de ejecución y posibilidad de funcionar en local, sin depender de servidores externos.

El sistema trabaja con cuatro modelos principales:

\begin{itemize}
	\item \textbf{TF-IDF con Regresión Logística.} Es el modelo base. Convierte el texto en números mediante \gls{tfidf} y después aplica un clasificador de Regresión Logística.
	
	\item \textbf{TF-IDF con SVD y Regresión Logística.} Añade una reducción de dimensionalidad mediante \gls{truncatedsvd} antes de clasificar.
	
	\item \textbf{TF-IDF con SVM.} Utiliza la representación \gls{tfidf} y la conecta directamente con una Máquina de Vectores de Soporte, o \gls{svm}, con kernel lineal.
	
	\item \textbf{TF-IDF con RNA.} Utiliza \gls{tfidf} como entrada de una red neuronal artificial sencilla, concretamente un Perceptrón Multicapa o \gls{mlp}.
\end{itemize}

Comparar estos cuatro modelos permite observar si realmente compensa utilizar técnicas más complejas o si un modelo más sencillo puede ofrecer resultados suficientemente buenos. Esta comparación es importante porque el objetivo no es construir el sistema más pesado posible, sino uno que funcione bien con recursos limitados.

\section{Clasificación supervisada en texto clínico}

La tarea principal del sistema se plantea como un problema de clasificación supervisada multiclase. Esto significa que el modelo aprende a partir de ejemplos ya etiquetados: cada texto de síntomas del conjunto de entrenamiento tiene asociada una especialidad médica correcta. A partir de esos ejemplos, el sistema aprende patrones que luego utiliza para clasificar nuevos textos escritos por otros pacientes~\cite{ibm_clasificacion, LopezUbeda2025}.

En este proyecto, las clases posibles son las especialidades médicas. Por ejemplo, un texto relacionado con dolor torácico podría estar asociado a Cardiología, mientras que una descripción sobre lesiones en la piel podría relacionarse con Dermatología. El objetivo del modelo es recibir una descripción nueva y elegir la especialidad más probable.

Un aspecto importante es que el sistema no debe limitarse a memorizar frases del entrenamiento. Dos pacientes pueden describir el mismo problema con palabras diferentes. Por ejemplo, una persona puede escribir ``me falta el aire'' y otra ``tengo sensación de ahogo''. Aunque las frases no sean iguales, el sistema debería aprender que ambas pueden estar relacionadas con un problema respiratorio. Por eso se utilizan técnicas de regularización y evaluación que ayudan a evitar el sobreajuste, es decir, que el modelo funcione muy bien con los ejemplos conocidos pero mal con casos nuevos~\cite{aws_sobreajuste}.

\section{Preprocesamiento del texto}

Antes de aplicar cualquier modelo, el texto escrito por el paciente debe prepararse. Las frases pueden incluir mayúsculas, signos de puntuación, errores ortográficos, palabras sin valor clínico o distintas formas de expresar una misma idea. El preprocesamiento intenta reducir esa variabilidad para que el modelo trabaje con una versión más limpia y homogénea del texto~\cite{snowflake_pln, ibm_stemming}.

En primer lugar, el texto se convierte a minúsculas y se eliminan caracteres que no aportan información útil. Se conservan las letras propias del español, como las vocales con tilde y la ñ, porque forman parte natural del idioma y pueden afectar al tratamiento posterior de las palabras.

Después se eliminan las palabras vacías o \textit{stop-words}. Son palabras muy frecuentes, como artículos, preposiciones o conjunciones, que aparecen en casi cualquier frase pero no ayudan demasiado a distinguir una especialidad médica de otra. Por ejemplo, palabras como ``el'', ``la'', ``de'' o ``y'' suelen aportar menos información que términos como ``fiebre'', ``dolor'', ``mareo'' o ``tos''.

Por último, se aplica \gls{stemming} con el algoritmo Snowball para español. Esta técnica reduce las palabras a una raíz aproximada. Por ejemplo, palabras parecidas pueden quedar representadas de una forma similar, lo que ayuda a agrupar variantes de un mismo término. A diferencia de la lematización, que intenta obtener la forma correcta de una palabra según el idioma, el \gls{stemming} es más simple y rápido~\cite{ibm_stemming}. Esta elección encaja con el enfoque del proyecto, ya que se busca un sistema ligero y fácil de ejecutar en local~\cite{snowflake_pln}.

El mismo preprocesamiento se aplica tanto durante el entrenamiento como cuando el paciente utiliza la aplicación. Esto evita que el modelo reciba textos tratados de forma distinta en cada fase.

\section{Representación del texto mediante TF-IDF}

Una vez limpiado el texto, es necesario convertirlo en números, porque los algoritmos de aprendizaje automático no trabajan directamente con frases. Para ello se utiliza \gls{tfidf}, una técnica clásica de representación de textos~\cite{wikipedia_tfidf}.

La idea de \gls{tfidf} es dar más importancia a las palabras que son relevantes dentro de un texto, pero que no aparecen en todos los documentos. Por ejemplo, si una palabra aparece en casi todos los textos, probablemente no ayude mucho a diferenciar especialidades. En cambio, términos más concretos como ``arritmia'', ``erupción'', ``migraña'' o ``fractura'' pueden aportar más información para clasificar el caso.

Esta representación tiene dos ventajas importantes para el proyecto. La primera es que es sencilla y rápida de calcular. La segunda es que resulta relativamente interpretable, porque cada valor del vector está relacionado con una palabra o conjunto de palabras del vocabulario. Esto permite entender mejor qué términos han influido en la decisión del modelo.

En la implementación se utiliza \texttt{TfidfVectorizer} de scikit-learn. Además, se limita el número máximo de términos mediante el parámetro \texttt{max\_features}. Esta limitación evita que el vocabulario crezca demasiado y ayuda a controlar tanto el consumo de memoria como el tiempo de entrenamiento.

En algunos modelos se utilizan solo palabras individuales, mientras que en otros también se incluyen combinaciones de dos palabras. Estas combinaciones, llamadas bigramas, pueden ser útiles en expresiones clínicas como ``dolor agudo'', ``insuficiencia renal'' o ``dolor torácico''. Sin embargo, también aumentan el tamaño del vocabulario, por lo que se utilizan principalmente en los modelos donde pueden aportar una mejora real.

La matriz generada por \gls{tfidf} suele ser dispersa, es decir, contiene muchos valores a cero. Esto ocurre porque cada paciente utiliza solo una pequeña parte de todas las palabras posibles del vocabulario. Esta característica permite almacenar la información de forma eficiente y hace que modelos como el SVM lineal puedan entrenarse directamente sobre esta representación.

\section{Reducción de dimensionalidad con TruncatedSVD}

El modelo basado en \gls{truncatedsvd} intenta simplificar la representación generada por \gls{tfidf}. Cuando el vocabulario es grande, cada texto queda representado por muchos valores, aunque la mayoría sean cero. Esto puede hacer que el modelo trabaje con una cantidad de información demasiado grande o repetitiva.

\gls{truncatedsvd} reduce ese espacio a un número menor de componentes, conservando la información más relevante. En este proyecto se utilizan 100 componentes, un valor que busca mantener suficiente información para clasificar los síntomas, pero reduciendo el coste de cálculo.

Esta técnica está relacionada con el Análisis Semántico Latente o \gls{lsa}, que intenta encontrar relaciones generales entre términos y documentos mediante técnicas de reducción de dimensionalidad~\cite{ibm_lsa}. En la práctica, puede ayudar a que textos con palabras distintas pero significado parecido queden más próximos entre sí. Por ejemplo, ``ahogo'' y ``disnea'' no son la misma palabra, pero pueden aparecer en contextos clínicos parecidos.

Después de aplicar \gls{truncatedsvd}, se utiliza de nuevo un clasificador de Regresión Logística. De esta forma se puede comparar el modelo TF-IDF básico con otro que usa el mismo clasificador, pero trabajando sobre una representación reducida. Así se observa si la reducción de dimensionalidad mejora realmente los resultados o si elimina información útil.

\section{Clasificación mediante SVM}

El modelo \gls{svm} se utiliza directamente sobre la representación \gls{tfidf}. En este caso no se aplica antes \gls{truncatedsvd}, porque las SVM se utilizan habitualmente en tareas de clasificación supervisada y pueden trabajar con representaciones numéricas de alta dimensionalidad~\cite{ibm_svm, ibm_clasificacion}.

De forma sencilla, una \gls{svm} intenta separar las clases buscando una frontera de decisión lo más clara posible. En este proyecto, esa frontera sirve para distinguir entre las distintas especialidades médicas. Si el texto de un paciente contiene términos más próximos a los ejemplos de una especialidad concreta, el modelo tenderá a asignarlo a esa especialidad.

Se utiliza un kernel lineal porque es adecuado para problemas de texto con muchas variables. Además, suele ser más eficiente que otros kernels más complejos. Esta decisión encaja con el objetivo del proyecto: obtener un buen rendimiento sin aumentar innecesariamente el coste computacional.

El modelo también se configura para devolver una estimación de probabilidad. Esto es importante porque la aplicación no solo muestra la especialidad predicha, sino también un porcentaje de confianza. En el caso de las SVM, esta probabilidad no aparece de forma directa, sino que se obtiene mediante un proceso de calibración posterior. Aunque esto añade algo de coste durante el entrenamiento, permite ofrecer una salida más comprensible para el usuario.

\section{Clasificación mediante RNA}

El cuarto modelo utiliza una \gls{rna}, concretamente un Perceptrón Multicapa o \gls{mlp}. Esta red neuronal recibe como entrada la representación \gls{tfidf} del texto y aprende a relacionarla con las especialidades médicas. Las redes neuronales artificiales se utilizan en aprendizaje automático para aprender relaciones entre datos de entrada y salidas esperadas a partir de ejemplos~\cite{ibm_redes_neuronales, cienciadedatos_redes_neuronales}.

La ventaja de una red neuronal es que puede capturar relaciones más complejas que un modelo lineal. Sin embargo, también suele requerir más tiempo de entrenamiento y más control para evitar que memorice demasiado los datos. Por eso, en este proyecto se utiliza una red sencilla y se limita el tamaño del vocabulario de entrada.

Además, se aplica parada temprana o \textit{early stopping}. Esta técnica detiene el entrenamiento cuando el modelo deja de mejorar sobre un conjunto de validación. Así se evita seguir entrenando sin obtener mejoras reales y se reduce el riesgo de sobreajuste~\cite{aws_sobreajuste}.

También se utiliza regularización mediante el parámetro \texttt{alpha}, que ayuda a controlar la complejidad de la red~\cite{interactive_chaos_reg}. En conjunto, estas decisiones permiten probar una técnica más flexible que los modelos lineales, pero manteniendo el sistema dentro de unos límites razonables de coste y simplicidad.

\section{Optimización y evaluación de los modelos}

Para que la comparación entre modelos sea justa, no basta con entrenarlos una vez y mirar el resultado. Algunos algoritmos dependen mucho de ciertos parámetros, como el valor de regularización en una \gls{svm}, el tamaño del vocabulario o la estructura de una red neuronal. Por eso, en los modelos más complejos se aplica \gls{gscv}, que prueba varias combinaciones de parámetros y selecciona la que obtiene mejores resultados. Este tipo de ajuste permite comparar configuraciones de forma ordenada en lugar de escoger los valores manualmente~\cite{ibm_clasificacion, cienciadedatos_hiperparametros}.

La búsqueda de hiperparámetros se aplica principalmente a los modelos \gls{svm} y \gls{mlp}, porque son los que tienen más configuraciones posibles. En el caso de la \gls{svm}, se prueban distintos valores de regularización y diferentes tamaños de vocabulario. En el caso de la \gls{rna}, se prueban distintas arquitecturas, funciones de activación y valores de regularización.

Estas combinaciones se evalúan mediante validación cruzada estratificada con cinco particiones. La estratificación es importante porque no todas las especialidades tienen el mismo número de ejemplos. Si se repartieran los datos de forma completamente aleatoria, algunas particiones podrían quedar con muy pocos ejemplos de ciertas clases. La validación estratificada ayuda a mantener una distribución más equilibrada entre las clases~\cite{ibm_clasificacion}.

Los modelos TF-IDF base y SVD se utilizan como modelos de referencia. Por eso se entrenan con una configuración más sencilla, sin una búsqueda exhaustiva de hiperparámetros. Su función principal es servir como punto de comparación frente a modelos más costosos como \gls{svm} y \gls{mlp}.

\section{Métricas de evaluación}

Para medir el rendimiento se utilizan varias métricas habituales en clasificación supervisada~\cite{ibm_clasificacion, LopezUbeda2025}.

La métrica principal es la \textit{accuracy}, que indica el porcentaje total de predicciones correctas. Es fácil de interpretar y permite comparar rápidamente los modelos. Sin embargo, tiene una limitación: si hay clases con muchos más ejemplos que otras, un modelo puede obtener una \textit{accuracy} alta aunque funcione mal en las especialidades menos representadas.

Por eso también se analizan la \textit{precision}, el \textit{recall} y el \textit{F1-macro}. La \textit{precision} indica cuántas de las predicciones realizadas para una especialidad son correctas. El \textit{recall} indica cuántos casos reales de una especialidad consigue detectar el modelo. El \textit{F1-macro} resume el equilibrio entre \textit{precision} y \textit{recall}, dando el mismo peso a todas las clases, aunque unas tengan más ejemplos que otras.

En un sistema de orientación médica, estas métricas son especialmente importantes. No basta con acertar muchas veces en general; también interesa saber si el modelo falla demasiado en especialidades con menos ejemplos o si tiende a confundir síntomas que podrían requerir una derivación concreta.

\section{Umbral de confianza y derivación a atención general}

La aplicación no se limita a mostrar una especialidad. También muestra un porcentaje de confianza asociado a la predicción. Si el modelo tiene una confianza alta, se muestra la especialidad sugerida. Si la confianza es baja, el sistema pide al paciente que añada más información sobre sus síntomas.

Este mecanismo evita que el sistema fuerce una respuesta cuando no tiene suficiente seguridad. Si después de varios intentos la confianza sigue siendo baja, el caso se deriva a atención general y los síntomas se guardan en un registro para que puedan revisarse posteriormente.

Este enfoque se basa en la clasificación con umbral de confianza. En lugar de aceptar siempre la clase más probable, el sistema descarta las predicciones que no superan un mínimo de seguridad. Este tipo de umbral permite convertir las puntuaciones del modelo en decisiones más prudentes, especialmente cuando la predicción no alcanza suficiente confianza~\cite{google_umbral_clasificacion}. En un contexto médico, esta decisión es prudente, porque una asignación incorrecta puede ser más problemática que reconocer que el sistema no tiene información suficiente.

\section{IA Frugal y Edge AI}

El proyecto se ha diseñado siguiendo una filosofía de IA Frugal. Esto significa que se busca resolver el problema con el menor consumo de recursos posible, evitando modelos innecesariamente grandes o dependencias externas cuando no son imprescindibles. Este enfoque conecta con la idea de reducir el impacto y el coste de los sistemas de inteligencia artificial, manteniendo solo los recursos necesarios para cumplir el objetivo del proyecto~\cite{greenai_es}.

En este caso, la IA Frugal se refleja en varias decisiones. Se utiliza \gls{tfidf} porque es una representación ligera y suficiente para una primera aproximación al problema. Se aplica \gls{truncatedsvd} para reducir el tamaño de los datos. Se emplean modelos clásicos como Regresión Logística y \gls{svm}, y la red neuronal se mantiene en una arquitectura sencilla.

Además, el sistema está pensado para funcionar en local, siguiendo el enfoque de Edge AI. Esto significa que la inferencia se realiza en el propio dispositivo, sin enviar los síntomas del paciente a un servidor externo. Esta característica mejora la privacidad, reduce la dependencia de conexión a internet y facilita que el sistema pueda utilizarse en entornos con recursos limitados~\cite{ibm_edge_ai}.

\section{Conclusiones del capítulo}

Los conceptos explicados en este capítulo justifican la arquitectura del sistema. El texto del paciente se limpia, se transforma mediante \gls{tfidf} y se clasifica utilizando distintos modelos. Cada modelo permite comprobar una hipótesis concreta: si basta con una solución sencilla, si la reducción de dimensionalidad ayuda, si una \gls{svm} mejora el resultado o si una \gls{rna} aporta alguna ventaja adicional.

La propuesta mantiene un equilibrio entre sencillez, interpretabilidad y rendimiento. Además, el uso de un umbral de confianza permite que el sistema actúe con prudencia cuando no tiene seguridad suficiente, derivando el caso a atención general en lugar de forzar una especialidad concreta.

Aun así, el sistema tiene limitaciones. Al basarse en representaciones de tipo bolsa de palabras, no siempre interpreta bien la negación, los matices del lenguaje o expresiones clínicas ambiguas. Por ejemplo, puede resultar difícil diferenciar correctamente entre ``tengo fiebre'' y ``no tengo fiebre'' si el modelo no captura bien el contexto. Estas limitaciones se asumen como parte del alcance del proyecto y quedan abiertas como posibles mejoras futuras.